\lesson{2}{от 11 мар 2025 12:45}{Продолжение}

\subsection{Обучение с учителем}

\begin{definition}[Supervised learning]
  \emph{Sypervised learning} -- наиболее распространненный случай. Каждый прецедент представляет собой пару "объект, ответ". Требуется найти функциональную зависимость ответов от описаний объектов и построить алгоритм, принимающий на входе описание объекта и выдающий на выходе ответ.

  Функционал качества обычно определяется как средняя ошибка ответов, выданных алгоритмом по всем объектам выборки.
\end{definition}

\begin{definition}[Задача классификации, класс]
  \emph{Задача классификации} отличается тем, что множество допустимых ответов конечно. Их называют метками классов. \emph{Класс} -- это множество объектов с данным значением метки.
\end{definition}

\begin{definition}[Задача регрессии]
  \emph{Задача регрессии} отличается тем, что допустимым ответом ялвяется действительное число или числовой вектор.
\end{definition}

\begin{definition}[Задача ранжирования]
  \emph{Задача ранжирования} отличается тем, что ответы надо получить сразу на множестве объектов, после чего отсортировать их по значениям ответов. Может сводиться к задачам классификации или регрессии. Часто применяется в информационном поиске и анализе текстов.
\end{definition}

\begin{definition}[Задача прогнозирования]
  \emph{Задача прогнозирования} отличается тем, что обхектами являеются отрезки временных рядов, обрывающиеся в тот момент, когда требуется сделать прогноз на будущее. Для решения задач пронозирования часто удается приспособить методы регрессии...
\end{definition}

\subsection{Обучение без учителя}

\begin{definition}[Обучение без учителя]
  \emph{Обучение без учителя}. В этом случае ответы не задаются, и требуется искать зависимости между объектами.
\end{definition}

\begin{definition}[Задача кластеризации]
  \emph{Задача кластеризации} заключается в том, чтобы сгруппировать объекты в кластеры, используя данные о попарном сходстве объектов. Функционалы качетсва могут определяться по-разному, например, как отношение средних межкластерных и внутрикластерных расстояний.
\end{definition}

\subsection{Кластерный анализ}

\begin{definition}[DBSCAN]
  \emph{DBSCAN} -- это алгоритм кластеризации, основанный на плотности.
\end{definition}
